% !TeX program = xelatex
\documentclass[a4paper,10pt]{article}

% ============================================================================
% សន្លឹកកិច្ចការសិស្ស ស-ខ្មៅ (Student Worksheet - Grade 8 3.1.3 INVP)
% បរិបទ៖ គ្មានស្លាយ LCD និងគ្មានម៉ាស៊ីនគិតលេខ — សិស្សបំពេញផ្ទាល់លើសន្លឹកកិច្ចការ ២ ទំព័រ A4
% ============================================================================
\usepackage[a4paper, left=0.85cm, right=0.85cm, top=0.70cm, bottom=0.95cm, footskip=0.48cm]{geometry}
\usepackage{amsmath, amssymb}
\usepackage{unicode-math}
\usepackage{fontspec}
\usepackage{array, colortbl}
\usepackage[table]{xcolor}
\usepackage{enumitem}
\usepackage{pifont}
\usepackage{fancyhdr}
\usepackage{lastpage}
\usepackage{tikz}
\usepackage[most]{tcolorbox}
\usetikzlibrary{calc, arrows.meta, decorations.pathreplacing, positioning, babel}

\makeatletter
\let\@expl@@@mark@update@singlecol@structures@@\relax
\makeatother

\providecommand{\NGSFindFonts}{%
  \IfFileExists{C:/Windows/Fonts/cambria.ttc}{%
    \def\NGSSysPath{C:/Windows/Fonts/}%
    \def\NGSUserPath{C:/Users/KIT/AppData/Local/Microsoft/Windows/Fonts/}%
  }{%
    \IfFileExists{./fonts/KhmerOSsiemreap.ttf}{\def\NGSSysPath{./fonts/}\def\NGSUserPath{./fonts/}}{%
    \IfFileExists{../fonts/KhmerOSsiemreap.ttf}{\def\NGSSysPath{../fonts/}\def\NGSUserPath{../fonts/}}{%
    \IfFileExists{../../fonts/KhmerOSsiemreap.ttf}{\def\NGSSysPath{../../fonts/}\def\NGSUserPath{../../fonts/}}{%
    \IfFileExists{../../../fonts/KhmerOSsiemreap.ttf}{\def\NGSSysPath{../../../fonts/}\def\NGSUserPath{../../../fonts/}}{%
    \IfFileExists{../../../../fonts/KhmerOSsiemreap.ttf}{\def\NGSSysPath{../../../../fonts/}\def\NGSUserPath{../../../../fonts/}}{%
    \IfFileExists{../../../../../fonts/KhmerOSsiemreap.ttf}{\def\NGSSysPath{../../../../../fonts/}\def\NGSUserPath{../../../../../fonts/}}{%
    \IfFileExists{../../../../../../fonts/KhmerOSsiemreap.ttf}{\def\NGSSysPath{../../../../../../fonts/}\def\NGSUserPath{../../../../../../fonts/}}{%
      \def\NGSSysPath{../../../../../../../fonts/}\def\NGSUserPath{../../../../../../../fonts/}%
    }}}}}}}%
  }%
}
\NGSFindFonts

\setmathfont[Path={\NGSSysPath}, FontIndex=1, Scale=1.02]{cambria.ttc}
\defaultfontfeatures{Ligatures=TeX}

\setmainfont[
  Path={\NGSUserPath}, Extension=.ttf, Script=Khmer,
  BoldFont={KhmerOSsiemreap}, BoldFeatures={FakeBold=1.8},
  ItalicFont={KhmerOSsiemreap}, ItalicFeatures={FakeSlant=0.22}
]{KhmerOSsiemreap}

\newfontfamily\muol[Path={\NGSUserPath}, Extension=.ttf, Script=Khmer]{KhmerOSmuollight}
\newfontfamily\TimesEn[Path={\NGSSysPath}, Extension=.ttf, BoldFont={timesbd}, ItalicFont={timesi}]{times}

\XeTeXlinebreaklocale "km"
\XeTeXlinebreakskip = 0pt plus 1pt minus 0.1pt

\definecolor{panelgray}{gray}{0.93}
\definecolor{headgray}{gray}{0.85}

\pagestyle{fancy}
\fancyhf{}
\renewcommand{\headrulewidth}{0pt}
\fancyfoot[L]{\fontsize{7.5pt}{9.5pt}\selectfont អនុវិទ្យាល័យព្រះក្រសាំង \en{|} គណិតវិទ្យាថ្នាក់ទី៨ \en{|} មេរៀនទី៣៖ ១.៣. ទំហំសមាមាត្រច្រាស (គ្រូ៖ ហុក ប៊ុនថែន)}
\fancyfoot[R]{\fontsize{7.8pt}{9.5pt}\selectfont ទំព័រ \en{\thepage/\pageref{LastPage}}}

\setlength{\parindent}{0pt}
\setlength{\parskip}{0pt}
\setlength{\tabcolsep}{3.8pt}

\newcommand{\en}[1]{{\TimesEn #1}}
\newcommand{\enbf}[1]{{\TimesEn\bfseries #1}}
\newcommand{\unit}[1]{\,\text{\en{#1}}}
\newcommand{\cnum}[1]{\mbox{\fontsize{10pt}{12pt}\selectfont\ding{\numexpr181+#1\relax}}}
\newcommand{\chkbox}{\tikz[baseline=-0.12ex]\draw[black, line width=0.8pt] (0,0) rectangle (0.24cm,0.24cm);}
\newcommand{\blank}[1][1.5cm]{\makebox[#1]{\dotfill}}
\newcommand{\ansbox}[1][0.85cm]{\,\tikz[baseline=-0.5ex]\draw[black, line width=0.65pt, fill=white] (0,0) rectangle (#1, 0.42cm);\,}
\newcommand{\secbar}[2]{%
  \begin{tcolorbox}[enhanced, sharp corners, colback=headgray, colframe=black, boxrule=0.9pt, boxsep=1pt, top=2pt, bottom=2pt, left=5pt, right=5pt, before skip=4pt, after skip=3pt]
    {\muol\fontsize{9.2pt}{13pt}\selectfont #1}\hfill{\fontsize{8pt}{11pt}\selectfont\textbf{#2}}
  \end{tcolorbox}%
}

\newcolumntype{L}[1]{>{\raggedright\arraybackslash}p{#1}}
\newcolumntype{C}[1]{>{\centering\arraybackslash}p{#1}}

\begin{document}
\fontsize{9.6pt}{14.4pt}\selectfont

% ============================================================================
% ទំព័រទី ១៖ E1 + E2 (POE 1 & POE 2) + E3 សំយោគនិយមន័យ និងធនាគារហេតុផល
% ============================================================================
\begin{tcolorbox}[enhanced, sharp corners, colback=white, colframe=black, boxrule=1.0pt, boxsep=2pt, top=3pt, bottom=3pt, left=6pt, right=6pt]
  \begin{center}
    {\muol\fontsize{10.4pt}{15pt}\selectfont អនុវិទ្យាល័យព្រះក្រសាំង \enbf{--} សន្លឹកកិច្ចការសិស្ស \enbf{(}មេរៀនទី៣៖ ១\enbf{.}៣\enbf{.} ទំហំសមាមាត្រច្រាស\enbf{)}}\\[3pt]
    {\fontsize{8.8pt}{12pt}\selectfont
    ឈ្មោះសិស្ស / ដៃគូ៖ \blank[5.2cm] \hspace{0.25cm}
    ថ្នាក់ទី៨៖ \blank[1.1cm] \hspace{0.25cm}
    ក្រុមទី៖ \blank[1.0cm] \hspace{0.25cm}
    កាលបរិច្ឆេទ៖ ២៧/១១/២០២៦}
  \end{center}
\end{tcolorbox}

% --- E1 + E2 ---
\secbar{ផ្នែកទី១៖ បេសកកម្មទស្សន៍ទាយ \enbf{(E1)} និងការរុករកជាដៃគូ \enbf{(E2 -- POE 1 \& POE 2)}}{រយៈពេល ១៤ នាទី}

\begin{tcolorbox}[enhanced, sharp corners, colback=panelgray, colframe=black, boxrule=0.75pt, boxsep=1.5pt, top=2.5pt, bottom=2.5pt, left=5pt, right=5pt]
  \textbf{បេសកកម្ម \en{E1} (គូសធីក \chkbox\ មុនពិសោធន៍)៖}
  (១) កម្មករ $35$ នាក់សង់ផ្ទះចប់ក្នុង $16$ ថ្ងៃ — បើនៅសល់កម្មករ $28$ នាក់ តើត្រូវប្រើពេល៖ \chkbox\ $12.8$ ថ្ងៃ \quad \chkbox\ $20$ ថ្ងៃ\\
  (២) ចតុកោណកែងក្រឡាផ្ទៃថេរ $36\unit{m}^2$ — តើបរិមាណមួយណាថេរគ្រប់ករណី៖ \chkbox\ ផលធៀប $\frac{y}{x}$ \quad \chkbox\ ផលគុណ $x\times y$
\end{tcolorbox}

\vspace{3pt}
\textbf{ក. សកម្មភាពទី ១ (\en{POE 1})៖ គ្រួសារចតុកោណកែងក្រឡាផ្ទៃថេរ $S = 36\unit{m}^2$ (សៀវភៅពុម្ព ទំ.៣០)}

{\renewcommand{\arraystretch}{1.48}
\begin{tabular}{|L{5.10cm}|C{2.05cm}|C{2.05cm}|C{2.05cm}|C{2.05cm}|C{2.05cm}|C{2.05cm}|}
  \hline
  \rowcolor{headgray}
  \textbf{បរិមាណនៃចតុកោណកែង} & \textbf{ទី ១} & \textbf{ទី ២} & \textbf{ទី ៣} & \textbf{ទី ៤} & \textbf{ទី ៥} & \textbf{ទី ៦} \\
  \hline
  ប្រវែងទទឹង $x$ ($\unit{m}$) & $1$ & $2$ & $3$ & $4$ & $6$ & $9$ \\
  \hline
  ប្រវែងបណ្ដោយ $y$ ($\unit{m}$) & $36$ & $18$ & $\blank[1.0cm]$ & $\blank[1.0cm]$ & $6$ & $\blank[1.0cm]$ \\
  \hline
  ផលគុណ $x \times y$ ($\unit{m}^2$) & $36$ & $\blank[1.0cm]$ & $\blank[1.0cm]$ & $\blank[1.0cm]$ & $\blank[1.0cm]$ & $\blank[1.0cm]$ \\
  \hline
\end{tabular}}

\vspace{3pt}
$\bullet$ \textbf{ពន្យល់ (\en{E})៖} ពេលទទឹង $x$ កើន $\times 2$ នោះបណ្ដោយ $y$ ថយ $\div \blank[0.8cm]$ ហើយផលគុណ $x\cdot y = \blank[1.0cm]$ ថេរជានិច្ច $\implies y = \dfrac{\blank[0.9cm]}{x}$ ។

\vspace{4pt}
\textbf{ខ. សកម្មភាពទី ២ (\en{POE 2})៖ ញែកដាច់ សមាមាត្រស្រប ទល់នឹង សមាមាត្រច្រាស ទល់នឹង អន្ទាក់ផលបូកថេរ}

\begin{tabular}{|L{9.20cm}|L{9.20cm}|}
  \hline
  \textbf{តារាង ក៖ ល្បឿន $v$ និងរយៈពេល $t$ សម្រាប់ចម្ងាយ $120\unit{km}$}\par\vspace{2pt}
  $\bullet$ $(v, t)$៖ $(20, 6),\; (30, 4),\; (40, 3),\; (60, 2)$\par\vspace{2pt}
  $\bullet$ គណនាផលគុណ $v\cdot t$៖ $20\!\times\! 6 = 120$ ; $30\!\times\! 4 = \blank[0.9cm]$ ; $40\!\times\! 3 = \blank[0.9cm]$\par\vspace{2pt}
  $\bullet$ សន្និដ្ឋាន៖ \chkbox\ សមាមាត្រច្រាស \quad \chkbox\ មិនសមាមាត្រច្រាស
  &
  \textbf{តារាង ខ៖ ចាយប្រាក់ $x$ និងប្រាក់នៅសល់ $y$ (មានប្រាក់ $20\$$)}\par\vspace{2pt}
  $\bullet$ $(x, y)$៖ $(2, 18),\; (5, 15),\; (10, 10)$ ($x$ កើន ហើយ $y$ ថយ!)\par\vspace{2pt}
  $\bullet$ គណនាផលគុណ $x\cdot y$៖ $2\!\times\! 18 = 36$ ; $5\!\times\! 15 = \blank[0.9cm]$ ; $10\!\times\! 10 = \blank[0.9cm]$\par\vspace{2pt}
  $\bullet$ សន្និដ្ឋាន៖ \chkbox\ សមាមាត្រច្រាស \quad \chkbox\ មិនសមាមាត្រច្រាស \\
  \hline
\end{tabular}

\vspace{3pt}
$\bullet$ \textbf{ពន្យល់ (\en{E})៖} ហេតុអ្វីក្នុងតារាង ខ ពេល $x$ កើន ហើយ $y$ ថយ តែមិនមែនសមាមាត្រច្រាស? ព្រោះផលគុណ $x\cdot y$ \blank[2.2cm] ។

% --- E3 ---
\secbar{ផ្នែកទី២៖ សំយោគក្បួនទូទៅ \enbf{(E3)} និងធនាគារហេតុផល \enbf{B\&W} \cnum{1}\enbf{--}\cnum{5}}{សៀវភៅពុម្ព ទំ.៣០}

\begin{tabular}{|L{9.20cm}|L{9.20cm}|}
  \hline
  \textbf{១. ក្បួនជាទូទៅ (ទំ.៣០)៖}\par\vspace{2pt}
  $\bullet$ ប្រសិនបើទំហំ $y$ \textbf{សមាមាត្រច្រាស}នឹងទំហំ $x$ គេបាន៖\par\vspace{3pt}
  {\centering $x \cdot y = \ansbox[0.9cm]$ \quad ឬ \quad $y = \dfrac{\ansbox[0.9cm]}{x}$ \quad ($a \neq 0$)\par}\vspace{3pt}
  $\bullet$ ចំពោះពីរគូតម្លៃ $(x_1, y_1)$ និង $(x_2, y_2)$៖ $\mathbf{x_1 y_1 = \blank[1.5cm]}$ ។
  &
  \textbf{២. ធនាគារហេតុផល \en{B\&W} (\cnum{1}--\cnum{5})៖}\par\vspace{2pt}
  \cnum{1} និយមន័យ៖ $y$ សមាមាត្រច្រាសនឹង $x \iff x\cdot y = a$ ថេរ ($y=\frac{a}{x}$) ។\par\vspace{2pt}
  \cnum{2} ពហុគុណច្រាស៖ បើ $x$ គុណនឹង $k$ នោះ $y$ ត្រូវចែកនឹង $k$ ។\par\vspace{2pt}
  \cnum{3} ផលគុណស្មើគ្នា៖ $x_1 y_1 = x_2 y_2$ (បរិមាណការងារសរុបថេរ) ។\par\vspace{2pt}
  \cnum{4} ប្រយ័ត្នអន្ទាក់៖ ហាមគុណខ្វែងតាមសមាមាត្រស្រប $\frac{x_1}{y_1}=\frac{x_2}{y_2}$! \par\vspace{2pt}
  \cnum{5} អន្ទាក់ដកថេរ៖ $x + y = c$ មិនមែនសមាមាត្រច្រាសទេ ។ \\
  \hline
\end{tabular}

\vspace{3pt}
\begin{tcolorbox}[enhanced, sharp corners, colback=panelgray, colframe=black, boxrule=0.65pt, boxsep=1.5pt, top=2pt, bottom=2pt, left=5pt, right=5pt]
  \textbf{ប្រយ័ត្នអន្ទាក់៖} (១) សមាមាត្រស្របប្រើ\textbf{ផលធៀបថេរ} $\frac{y}{x}=a$ រីឯសមាមាត្រច្រាសប្រើ\textbf{ផលគុណថេរ} $xy=a$! \quad (២) កម្មករថយចុះ ចំនួនថ្ងៃត្រូវកើនឡើង!
\end{tcolorbox}

\vspace{3pt}
\textbf{៣. លំហាត់ត្រួតពិនិត្យរហ័ស៖ គណនាមេគុណសមាមាត្រច្រាស $a = xy$ និងតម្លៃមិនស្គាល់}

\vspace{2pt}
{\renewcommand{\arraystretch}{1.62}
\begin{tabular}{|C{4.40cm}|C{4.40cm}|C{4.40cm}|C{4.40cm}|}
  \hline
  \rowcolor{panelgray}
  $xy = 24,\; x=3 \Rightarrow y = \blank[0.9cm]$ & $x=6, y=8 \Rightarrow a = \blank[1.0cm]$ & $y=\dfrac{48}{x},\; x=12 \Rightarrow y = \blank[0.8cm]$ & $y=\dfrac{48}{x},\; y=16 \Rightarrow x = \blank[0.8cm]$ \\
  \hline
  $xy = -30,\; x=5 \Rightarrow y = \blank[0.9cm]$ & $x_1=4, y_1=9, x_2=6 \Rightarrow y_2=\blank[0.8cm]$ & $y=\dfrac{100}{x},\; x=25 \Rightarrow y=\blank[0.8cm]$ & $x+y=10$៖ ច្រាស? \blank[0.9cm] \\
  \hline
  $xy = 72,\; x=8 \Rightarrow y = \blank[0.9cm]$ & $x=5, y=14 \Rightarrow a = \blank[1.0cm]$ & $y=\dfrac{90}{x},\; x=15 \Rightarrow y = \blank[0.8cm]$ & $y=\dfrac{90}{x},\; y=18 \Rightarrow x = \blank[0.8cm]$ \\
  \hline
  $x_1=10, y_1=6, x_2=12 \Rightarrow y_2=\blank[0.7cm]$ & $x_1=15, y_1=4, y_2=10 \Rightarrow x_2=\blank[0.7cm]$ & $y=\dfrac{-36}{x},\; x=9 \Rightarrow y=\blank[0.8cm]$ & $y = 4x$៖ ច្រាស? \blank[0.9cm] \\
  \hline
  $xy = 60,\; x=15 \Rightarrow y = \blank[0.9cm]$ & $x=9, y=7 \Rightarrow a = \blank[1.0cm]$ & $y=\dfrac{75}{x},\; x=25 \Rightarrow y=\blank[0.8cm]$ & $xy = 12$៖ ច្រាស? \blank[0.9cm] \\
  \hline
\end{tabular}}

\newpage

% ============================================================================
% ទំព័រទី ២៖ E4 លំហាត់គំរូ + E5 ប្រតិបត្តិ + វិនិច្ឆ័យ A-D + Exit Check + កិច្ចការផ្ទះ
% ============================================================================
\secbar{ផ្នែកទី៣៖ ពង្រីកចំណេះដឹងតាមច្រកផលគុណថេរ \enbf{(E4 --} គំរូ ទំ\enbf{.}៣០\enbf{--}៣១\enbf{)}}{រយៈពេល ១០ នាទី}

\textbf{១. លំហាត់គំរូ (ទំ.៣០–៣១)៖} កម្មករ $35$ នាក់សង់ផ្ទះមួយហើយក្នុងរយៈពេល $16$ ថ្ងៃ ។ បើកម្មករតែ $28$ នាក់ធ្វើការក្នុងល្បឿនដូចគ្នា តើគេត្រូវការរយៈពេលប៉ុន្មានថ្ងៃដើម្បីសង់ផ្ទះនោះឱ្យហើយ?

{\renewcommand{\arraystretch}{1.48}
\begin{tabular}{|L{4.60cm}|L{9.00cm}|C{4.20cm}|}
  \hline
  \rowcolor{headgray}
  \textbf{ជំហានស៊ើបអង្កេត} & \textbf{ការវិភាគ និងការគណនាដោយដៃ} & \textbf{ចម្លើយ និងហេតុផល} \\
  \hline
  \textbf{១. ពិនិត្យទិសដៅទំនងហេតុផល} & កម្មករថយចុះពី $35 \to 28$ នាក់ $\implies$ ចំនួនថ្ងៃ $x$ ត្រូវ៖ \chkbox\ តិចជាង $16$ \quad \chkbox\ ច្រើនជាង $16$ & សមាមាត្រច្រាស (\cnum{2}) \\
  \hline
  \textbf{២. សរសេរសមីការផលគុណថេរ} & បរិមាណការងារសរុបថេរ៖ $28 \times x = 35 \times \blank[1.0cm]$ & ហាមគុណខ្វែងស្រប! (\cnum{3}, \cnum{4}) \\
  \hline
  \textbf{៣. សម្រួលប្រភាគដោយដៃ} & $x = \dfrac{35 \times 16}{28} = \dfrac{5 \times 16}{\blank[0.8cm]} = 5 \times \blank[0.8cm] = \blank[1.2cm]$ & ត្រូវការ $\blank[1.0cm]$ ថ្ងៃ (\ding{51}) \\
  \hline
\end{tabular}}

\vspace{4pt}
\textbf{២. ពិនិត្យតារាង និងសរសេរសមីការ (លំនាំ \#6ឃ និង \#8 ទំ.៤០–៤១)៖}

{\renewcommand{\arraystretch}{1.44}
\begin{tabular}{|L{8.90cm}|L{8.90cm}|}
  \hline
  \textbf{ក.} $(x,y)$៖ $(2, 30),\; (3, 20),\; (5, 12),\; (6, 10)$\par\vspace{2pt}
  ផលគុណ $xy = \blank[1.1cm]$ ថេរ $\implies$ សមីការ៖ $y = \blank[1.6cm]$ (\cnum{1})
  &
  \textbf{ខ.} បើ $y$ សមាមាត្រច្រាសនឹង $x$ ហើយពេល $x = 8$ គេបាន $y = 9$៖\par\vspace{2pt}
  មេគុណ $a = 8\times 9 = \blank[1.0cm] \implies y = \blank[1.3cm]$ ; ពេល $x=12 \Rightarrow y=\blank[0.9cm]$ \\
  \hline
\end{tabular}}

\vspace{4pt}
$\bullet$ \textbf{សំណួរឆ្លុះបញ្ចាំង \en{(EIA Looking-Back)}៖} ហេតុអ្វីចម្លើយ $x = 12.8$ ថ្ងៃសម្រាប់កម្មករ $28$ នាក់ ដឹងថាខុសភ្លាមដោយមិនបាច់គិតលេខវែង? \blank[2.5cm]

% --- E5 ---
\secbar{ផ្នែកទី៤៖ វាយតម្លៃលទ្ធផលសិក្សា \enbf{(E5 --} ប្រតិបត្តិ ទំ\enbf{.}៣១ \enbf{+} វិនិច្ឆ័យ \enbf{A--D +} សំបុត្រចេញ\enbf{)}}{រយៈពេល ១០ នាទី}

\textbf{ក. លំហាត់ប្រតិបត្តិ (សៀវភៅពុម្ព ទំ.៣១ លើ — ធ្វើជាដៃគូ រួចប្ដូរសន្លឹកផ្ទៀង \en{Pause-Review})៖}

{\renewcommand{\arraystretch}{1.48}
\begin{tabular}{|L{8.90cm}|L{8.90cm}|}
  \hline
  \textbf{១) ប្រតិបត្តិ (ទំ.៣១ — ក្បាលរ៉ូប៊ីណេទឹក)៖} រ៉ូប៊ីណេដូចគ្នា $4$ បង្ហូរទឹកពេញអាងក្នុងរយៈពេល $70\unit{mn}$ ។ បើបើករ៉ូប៊ីណេ $7$ រួមគ្នា៖\par\vspace{2pt}
  $\bullet$ សមីការផលគុណថេរ៖ $7 \times t = 4 \times \blank[1.1cm] = \blank[1.3cm]$\par\vspace{2pt}
  $\bullet$ រយៈពេល៖ $t = \dfrac{\blank[1.2cm]}{7} = \blank[1.4cm]\unit{mn}$ (\cnum{1}, \cnum{3})
  &
  \textbf{២) ពង្រឹង (ល្បឿន និងរយៈពេលធ្វើដំណើរ)៖} រថយន្តបើកល្បឿន $60\unit{km/h}$ ប្រើពេល $5\unit{h}$ ដល់គោលដៅ ។ បើបើកល្បឿន $75\unit{km/h}$៖\par\vspace{2pt}
  $\bullet$ ចម្ងាយផ្លូវសរុបថេរ៖ $75 \times t = 60 \times 5 = \blank[1.3cm]\unit{km}$\par\vspace{2pt}
  $\bullet$ រយៈពេល៖ $t = \dfrac{300}{75} = \blank[1.4cm]\unit{h}$ (\cnum{3}) \\
  \hline
\end{tabular}}

\vspace{5pt}
\textbf{ខ. គូសវង់ជុំវិញជម្រើសត្រឹមត្រូវតែ ១ គត់ ក្នុងចំណោម \en{A, B, C, D} រួចសរសេរលេខហេតុផល \cnum{1}--\cnum{5} នៃជម្រើសខុស៖}

{\renewcommand{\arraystretch}{1.35}
\begin{tabular}{|C{4.40cm}|C{4.40cm}|C{4.40cm}|C{4.40cm}|}
  \hline
  \cellcolor{panelgray}\textbf{ជម្រើស \en{A}} & \cellcolor{panelgray}\textbf{ជម្រើស \en{B}} & \cellcolor{panelgray}\textbf{ជម្រើស \en{C}} & \cellcolor{panelgray}\textbf{ជម្រើស \en{D}} \\
  \hline
  រ៉ូប៊ីណេ $7$៖ $t = \frac{280}{7} = 40\unit{mn}$ & រ៉ូប៊ីណេ $7$៖ $t = 122.5\unit{mn}$ & $x + y = 20$ សមាមាត្រច្រាស & កម្មករ $28$ នាក់ សង់ $12.8$ ថ្ងៃ \\[3pt]
  \chkbox\ ត្រូវ \quad \chkbox\ ខុស (\blank[0.7cm]) & \chkbox\ ត្រូវ \quad \chkbox\ ខុស (\blank[0.7cm]) & \chkbox\ ត្រូវ \quad \chkbox\ ខុស (\blank[0.7cm]) & \chkbox\ ត្រូវ \quad \chkbox\ ខុស (\blank[0.7cm]) \\
  \hline
\end{tabular}}

\vspace{5pt}
\begin{tcolorbox}[enhanced, sharp corners, colback=white, colframe=black, boxrule=0.85pt, boxsep=2.5pt, top=3.5pt, bottom=3.5pt, left=5pt, right=5pt]
  {\muol\fontsize{8.8pt}{12.5pt}\selectfont សំបុត្រចេញរហ័ស \enbf{(Exit Check --} ៦០ វិនាទី\enbf{)}៖}\\[2pt]
  បើ $y$ សមាមាត្រច្រាសនឹង $x$ ហើយពេល $x=4$ គេបាន $y=15$៖
  (១) $a = xy = \blank[1.1cm]$ \;\en{|}\;
  (២) សមីការ $y = \blank[1.3cm]$ \;\en{|}\;
  (៣) ពេល $x=12 \Rightarrow y = \blank[1.1cm]$
\end{tcolorbox}

\vspace{5pt}
\begin{tcolorbox}[enhanced, sharp corners, colback=panelgray, colframe=black, boxrule=0.85pt, boxsep=2.5pt, top=3.5pt, bottom=3.5pt, left=5pt, right=5pt]
  {\muol\fontsize{8.8pt}{12.5pt}\selectfont កិច្ចការផ្ទះ \enbf{(}សៀវភៅពុម្ព ទំព័រ ៤០\enbf{--}៤១ លំហាត់ \enbf{\#6(}ឃ\enbf{,}ង\enbf{), \#8, \#9} និង \enbf{\#13)}៖}\\[2pt]
  $\bullet$ \textbf{លំហាត់ \#6(ឃ,ង), \#8, \#9៖} ពិនិត្យតារាង និងគណនាតម្លៃក្នុងទំហំសមាមាត្រច្រាស ។\\
  $\bullet$ \textbf{លំហាត់ \#13 (ទំ.៤១)៖} ដោះស្រាយចំណោទសមាមាត្រច្រាស និងអានមុន \textbf{២.១–២.២ ភាគរយ និងប្រាក់ចំណេញ-ខាត (ទំ.៣១–៣៣)} ។
\end{tcolorbox}

\end{document}
